\chapter{Boolean Analysis — Uniform Transfer}
%%% generated by the blueprint pipeline from TCSlib.BooleanAnalysis.GaussianSpace.Noise.UniformTransfer
\section{Overview}

Declarations of the Lean module \texttt{TCSlib.BooleanAnalysis.GaussianSpace.Noise.UniformTransfer}.

\section{Declarations}

\begin{definition}[Paired Gaussian torus map]
\lean{BooleanAnalysis.GaussianSpace.gaussianPairTorusMap}
\label{BooleanAnalysis.GaussianSpace.gaussianPairTorusMap}
\leanok
\uses{BooleanAnalysis.GaussianSpace.gaussianTorusMap}
For a natural number $n$, define a map from pairs $(z,w)$ of vectors in $\bbr^n$ to the
torus indexed by $(\mathrm{Fin}(n)\times\{0,1\})\sqcup(\mathrm{Fin}(n)\times\{0,1\})$ by
applying the Gaussian torus map to $z$ on the first summand and to $w$ on the second.
\end{definition}

\begin{lemma}[Continuity of the Gaussian pair torus map]
\lean{BooleanAnalysis.GaussianSpace.continuous_gaussianPairTorusMap}
\label{BooleanAnalysis.GaussianSpace.continuous_gaussianPairTorusMap}
\leanok
\uses{BooleanAnalysis.GaussianSpace.gaussianPairTorusMap, BooleanAnalysis.GaussianSpace.gaussianTorusMap}
\difficulty{4}
For each natural number $n$, the map from pairs $(x,y)$ of vectors in $\mathbb R^n$ to
the product torus indexed by the disjoint union of two copies of
$\{1,\ldots,n\}\times\{0,1\}$, obtained by applying the Gaussian torus map to each
vector, is continuous.
\end{lemma}

\begin{lemma}[Fourier character of the paired torus map]
\lean{BooleanAnalysis.GaussianSpace.mFourier_gaussianPairTorusMap}
\label{BooleanAnalysis.GaussianSpace.mFourier_gaussianPairTorusMap}
\leanok
\uses{BooleanAnalysis.GaussianSpace.gaussianPairTorusMap, BooleanAnalysis.GaussianSpace.gaussianTorusFrequency, BooleanAnalysis.GaussianSpace.gaussianTorusMap, BooleanAnalysis.GaussianSpace.mFourier_gaussianTorusMap}
\difficulty{3}
Let $n\ge 0$, let
$m:(\{1,\ldots,n\}\times\{0,1\})\sqcup(\{1,\ldots,n\}\times\{0,1\})\to\mathbb Z$, and
let $z=(x,y)\in\mathbb R^n\times\mathbb R^n$. The Fourier character of the paired torus
image of $z$ is
$\exp\!\left(i\sum_{j=1}^n\bigl[2\pi(m_{1,j,0}+\sqrt2\,m_{1,j,1})x_j+2\pi(m_{2,j,0}+\sqrt2\,m_{2,j,1})y_j\bigr]\right)$,
where $m_{1,j,b}$ and $m_{2,j,b}$ are the values of $m$ on the first and second copies,
respectively.
\end{lemma}

\begin{lemma}[Norm bound for Boolean noise averaging]
\lean{BooleanAnalysis.GaussianSpace.norm_booleanNoiseExpectation_le}
\label{BooleanAnalysis.GaussianSpace.norm_booleanNoiseExpectation_le}
\leanok
\uses{BooleanAnalysis.BoolCube, BooleanAnalysis.Hypercontractivity.noiseKernel, BooleanAnalysis.Hypercontractivity.noiseKernel_nonneg, BooleanAnalysis.Hypercontractivity.noiseKernel_sum_right, BooleanAnalysis.uniformWeight}
\difficulty{5}
Let $N$ be a natural number, let $0\leq\rho\leq1$, and let
$F:\{0,1\}^N\times\{0,1\}^N\to\mathbb C$ satisfy $\|F(x,y)\|\leq C$ for every $x,y$.
Then the norm of the sum of $F(x,y)$ weighted by the noise kernel and the uniform weight
obeys
\[
\left\|2^{-N}\sum_{x,y\in\{0,1\}^N}K_\rho(x,y)F(x,y)\right\|\leq C.
\]
\end{lemma}

\begin{definition}[Boolean block expectation functional]
\lean{BooleanAnalysis.GaussianSpace.booleanBlockExpectation}
\label{BooleanAnalysis.GaussianSpace.booleanBlockExpectation}
\leanok
\uses{BooleanAnalysis.BoolCube, BooleanAnalysis.GaussianSpace.gaussianPairTorusMap, BooleanAnalysis.GaussianSpace.normalizedCubeBlocks, BooleanAnalysis.Hypercontractivity.noiseKernel, BooleanAnalysis.uniformWeight}
For natural numbers $n,m$ and a real parameter $\rho$, define a continuous
complex-linear functional on the continuous complex-valued functions on the paired
torus. It is the sum over pairs of Boolean-cube points of the noise-kernel-weighted
evaluations at the corresponding pair of normalized block vectors, multiplied by the
uniform weight for a cube of dimension $nm$.
\end{definition}

\begin{lemma}[Boolean block expectation formula]
\lean{BooleanAnalysis.GaussianSpace.booleanBlockExpectation_apply}
\label{BooleanAnalysis.GaussianSpace.booleanBlockExpectation_apply}
\leanok
\uses{BooleanAnalysis.BoolCube, BooleanAnalysis.GaussianSpace.booleanBlockExpectation, BooleanAnalysis.GaussianSpace.gaussianPairTorusMap, BooleanAnalysis.GaussianSpace.normalizedCubeBlocks, BooleanAnalysis.Hypercontractivity.noiseKernel, BooleanAnalysis.uniformWeight}
\difficulty{1}
Let $n,m$ be natural numbers, $\rho\in\mathbb R$, and let $F$ be a continuous
complex-valued function on the torus indexed by
$(\mathrm{Fin}\,n\times\{0,1\})\sqcup(\mathrm{Fin}\,n\times\{0,1\})$. The Boolean block
expectation of $F$ is
\[
2^{-nm}\sum_{x,y\in\{0,1\}^{nm}}K_\rho(x,y)\,F\bigl(G(x,y)\bigr),
\]
where $K_\rho$ is the noise kernel and $G(x,y)$ is the pair of torus images obtained
from the normalized blocks of $x$ and $y$.
\end{lemma}

\begin{lemma}[Norm bound for the Boolean block expectation]
\lean{BooleanAnalysis.GaussianSpace.norm_booleanBlockExpectation_le}
\label{BooleanAnalysis.GaussianSpace.norm_booleanBlockExpectation_le}
\leanok
\uses{BooleanAnalysis.GaussianSpace.booleanBlockExpectation, BooleanAnalysis.GaussianSpace.booleanBlockExpectation_apply, BooleanAnalysis.GaussianSpace.norm_booleanNoiseExpectation_le}
\difficulty{1}
For natural numbers $n,m$, a real number $\rho$ with $0\le\rho\le1$, and a continuous
complex-valued function $F$ on the torus indexed by
$(\{1,\ldots,n\}\times\{0,1\})\sqcup(\{1,\ldots,n\}\times\{0,1\})$, the norm of its
Boolean block expectation, formed using blocks of size $m$ and correlation parameter
$\rho$, is at most the supremum norm of $F$.
\end{lemma}

\begin{definition}[Integration on a compact space]
\lean{BooleanAnalysis.GaussianSpace.compactIntegralCLM}
\label{BooleanAnalysis.GaussianSpace.compactIntegralCLM}
\leanok
Let $X$ be a compact topological space equipped with a measurable structure whose
measurable sets include the Borel sets, and let $\mu$ be a finite measure on $X$.
Integration against $\mu$ defines a continuous complex-linear functional on the
continuous complex-valued functions on $X$, sending $F$ to $\int_X F\,d\mu$.
\end{definition}

\begin{lemma}[Compact integral operator evaluation]
\lean{BooleanAnalysis.GaussianSpace.compactIntegralCLM_apply}
\label{BooleanAnalysis.GaussianSpace.compactIntegralCLM_apply}
\leanok
\uses{BooleanAnalysis.GaussianSpace.compactIntegralCLM}
\difficulty{1}
Let $X$ be a compact topological space equipped with a measurable structure whose
measurable sets include its Borel sets, let $\mu$ be a finite measure on $X$, and let
$F:X\to\mathbb C$ be continuous. The continuous linear functional defined by integration
against $\mu$ evaluates at $F$ to $\int_X F\,d\mu$.
\end{lemma}

\begin{lemma}[Convergence of uniformly bounded linear functionals]
\lean{BooleanAnalysis.GaussianSpace.tendsto_clm_of_dense_span}
\label{BooleanAnalysis.GaussianSpace.tendsto_clm_of_dense_span}
\leanok
\difficulty{5}
Let $E$ be a complex normed vector space, let $(L_m)_{m\in\mathbb N}$ and $M$ be
continuous complex-linear functionals from $E$ to $\mathbb C$, and let $C\in\mathbb R$.
Suppose that $\|L_m(y)\|\le C\|y\|$ for every $m$ and $y\in E$, and that the complex
linear span of a set $S\subseteq E$ is dense in $E$. If $L_m(y)\to M(y)$ for every $y\in
S$, then $L_m(x)\to M(x)$ for every $x\in E$.
\end{lemma}

\begin{lemma}[Convergence of Boolean block expectations]
\lean{BooleanAnalysis.GaussianSpace.tendsto_booleanBlockExpectation}
\label{BooleanAnalysis.GaussianSpace.tendsto_booleanBlockExpectation}
\leanok
\uses{BooleanAnalysis.GaussianSpace.booleanBlockExpectation, BooleanAnalysis.GaussianSpace.booleanBlockExpectation_apply, BooleanAnalysis.GaussianSpace.compactIntegralCLM, BooleanAnalysis.GaussianSpace.compactIntegralCLM_apply, BooleanAnalysis.GaussianSpace.continuous_gaussianPairTorusMap, BooleanAnalysis.GaussianSpace.correlatedGaussianMeasure, BooleanAnalysis.GaussianSpace.gaussianPairTorusMap, BooleanAnalysis.GaussianSpace.mFourier_gaussianPairTorusMap, BooleanAnalysis.GaussianSpace.norm_booleanBlockExpectation_le, BooleanAnalysis.GaussianSpace.tendsto_clm_of_dense_span, BooleanAnalysis.GaussianSpace.tendsto_sum_noiseKernel_normalizedCubeBlocks_phase}
\difficulty{4}
Let $n\ge0$, let $0\le\rho\le1$, and let $F$ be a continuous complex-valued function on
the torus $(\mathbb R/\mathbb
Z)^{(\{1,\ldots,n\}\times\{0,1\})\sqcup(\{1,\ldots,n\}\times\{0,1\})}$. For a Boolean
vector $x\in\{0,1\}^{nm}$, form $X_m(x)_i=m^{-1/2}\sum_{j=1}^m\sigma(x_{i,j})$, where
$\sigma(0)=1$ and $\sigma(1)=-1$; define $T(X_m(x))_{i,0}=X_m(x)_i\pmod{\mathbb Z}$ and
$T(X_m(x))_{i,1}=\sqrt2X_m(x)_i\pmod{\mathbb Z}$, and use the corresponding map on pairs
coordinatewise. The uniform average over $x,y\in\{0,1\}^{nm}$, weighted by
$2^{-nm}\prod_{i,j}(1+\rho\sigma(x_{i,j})\sigma(y_{i,j}))/2$, of
$F(T(X_m(x)),T(X_m(y)))$ converges as $m\to\infty$ to the integral of $F(T(g),T(h))$
under the law of $(G,\rho G+\sqrt{1-\rho^2}H)$, where $G,H$ are independent standard
Gaussian vectors in $\mathbb R^n$.
\end{lemma}

\begin{lemma}[Real Boolean block expectation convergence]
\lean{BooleanAnalysis.GaussianSpace.tendsto_booleanBlockExpectation_real}
\label{BooleanAnalysis.GaussianSpace.tendsto_booleanBlockExpectation_real}
\leanok
\uses{BooleanAnalysis.BoolCube, BooleanAnalysis.GaussianSpace.booleanBlockExpectation_apply, BooleanAnalysis.GaussianSpace.correlatedGaussianMeasure, BooleanAnalysis.GaussianSpace.gaussianPairTorusMap, BooleanAnalysis.GaussianSpace.normalizedCubeBlocks, BooleanAnalysis.GaussianSpace.tendsto_booleanBlockExpectation, BooleanAnalysis.Hypercontractivity.noiseKernel, BooleanAnalysis.uniformWeight}
\difficulty{3}
Let $n\ge0$, let $0\le\rho\le1$, and let $F$ be a continuous real-valued function on the
torus $(\mathbb R/\mathbb
Z)^{(\{1,\ldots,n\}\times\{0,1\})\sqcup(\{1,\ldots,n\}\times\{0,1\})}$. For a Boolean
vector $x\in\{0,1\}^{nm}$, set $X_m(x)_i=m^{-1/2}\sum_{j=1}^m\sigma(x_{i,j})$, with
$\sigma(0)=1$ and $\sigma(1)=-1$, and map each coordinate $u$ to
$(u,\sqrt2u)\pmod{\mathbb Z}$. The weighted average of $F$ evaluated on the two mapped
block vectors, with weight $2^{-nm}\prod_{i,j}(1+\rho\sigma(x_{i,j})\sigma(y_{i,j}))/2$
for each pair $x,y\in\{0,1\}^{nm}$, converges as $m\to\infty$ to the integral of the
same function under the law of $(G,\rho G+\sqrt{1-\rho^2}H)$, where $G,H$ are
independent standard Gaussian vectors in $\mathbb R^n$.
\end{lemma}
